\chapter{La prima ottimizzazione, contro una risposta nota}
\label{ch:firstopt}

Il capitolo precedente ha dato al motore un punteggio e una bussola.
Questo chiude l'anello per la prima volta: una ricerca che propone
progetti, li valuta con la marcia e prosegue finché non riesce più a
migliorare. Il punto non è la risposta ma la verifica --- il
problema è posto così piccolo che la sua soluzione è nota in anticipo
sulla carta, così la macchina può esservi confrontata. Un anello che
atterra su una risposta nota pur non sapendone nulla è degno di
fiducia su problemi la cui risposta non è nota.

\section{La domanda}

L'anello si chiude? Proporre, valutare, cercare --- e atterrare dove
la teoria dice che sta l'ottimo?

\section{Come funziona, per esteso}

La prima ottimizzazione è deliberatamente unidimensionale, così che
la sua risposta esista in forma chiusa (una formula esatta, in
contrapposizione a un numero ottenuto per iterazione). L'unica
manopola è il \emph{rapporto di espansione} $\varepsilon = A_e/A_t$
--- l'area di uscita divisa per l'area di gola, l'unico numero che
dice quanto l'ugello lascia espandere il gas
(Capitolo~\ref{ch:nozzle}). Tutto il resto è tenuto fisso.

La risposta nota viene dal lemma dell'area di uscita. Si derivi la
spinta~\eqref{eq:route1} rispetto all'area di uscita di un ugello
ideale (a uscita uniforme): allargare l'uscita di $\mathrm{d}A_e$
aggiunge un anello di parete su cui la pressione locale di uscita
$p_e$ spinge in avanti, mentre la pressione ambiente $p_a$ spinge
all'indietro sullo stesso anello, quindi
\begin{equation}
\frac{\mathrm{d}J}{\mathrm{d}A_e} \;=\; p_e - p_a
\qquad\Longrightarrow\qquad
\text{ottimo esattamente dove } p_e(\varepsilon^{*}) = p_a .
\label{eq:arealemma}
\end{equation}
Un ugello il cui gas è ancora sopra l'ambiente all'uscita
($p_e > p_a$) si dice \emph{sotto-espanso}: non ha finito di
espandersi e conviene allungarlo. Una volta che il gas è stato
espanso sotto l'ambiente ($p_e < p_a$) l'ugello è
\emph{sovra-espanso}: l'ambiente circostante spinge ora all'indietro
più forte di quanto il gas spinga in avanti su ogni nuovo anello, e
ogni anello aggiuntivo costa spinta. L'ottimo è il punto di
attraversamento. Porre una pressione ambiente fissa dunque
$\varepsilon^{*}$ analiticamente --- la marcia deve solo essere
invertita per trovare il rapporto di espansione a cui la sua
pressione di uscita eguaglia quell'ambiente.

La ricerca in sé non sa nulla di tutto questo. È una \emph{ricerca a
sezione aurea}: un metodo senza derivate che mantiene un intervallo
$[a,b]$ noto per contenere il massimo, valuta l'obiettivo in due
punti interni collocati secondo la sezione aurea, scarta l'estremo
che non può contenere il massimo e ripete finché l'intervallo non è
più stretto di una tolleranza fissata. Ogni restringimento riusa una
delle due valutazioni precedenti, ed è questo che la sezione aurea
fa guadagnare. I suoi unici ingressi sono chiamate alla marcia e
letture di $J$ --- nessun gradiente, nessuna conoscenza del lemma.
Ogni valutazione della marcia è memorizzata su disco con chiave
$\varepsilon$, così un'esecuzione interrotta riprende a costo nullo,
e il paesaggio della Figura~\ref{fig:jeps} è disegnato dalle
valutazioni della ricerca stessa anziché da una spazzata separata.

La suite chiude poi tre anelli indipendenti attorno alla risposta.
Il rapporto di espansione trovato dalla ricerca alla cieca deve
coincidere con la forma chiusa. La pressione di uscita lì deve
eguagliare l'ambiente posto, entro una banda derivata dalla stima
d'errore a due risoluzioni della marcia stessa più l'ampiezza
dell'intervallo su cui la ricerca si è fermata. E il misuratore di
spigolo della Sezione~\ref{sec:corner} --- costruito due passi
prima, mai consultato dalla ricerca --- deve leggere zero in modo
indipendente sul progetto trovato.

\section{Che cosa è stato costruito}

\carrier{validation/a1\_driver\_eps.py}: il driver a sezione aurea,
la cache su disco, le tre verifiche di chiusura e un rigettatore ---
un obiettivo deliberatamente corrotto, con il termine
pressione--area $(p_e-p_a)A_e$ cancellato, che deve allora correre
monotonamente fino al bordo dell'intervallo invece di trovare un
ottimo interno, perché senza quel termine nulla penalizza mai la
sovra-espansione.

\section{Risultati di riferimento}

\textbf{4/4 PASS} (231~s, 13 valutazioni della marcia). La ricerca
alla cieca atterra a $\varepsilon = 5.1551$ contro il valore
analitico $5.1506$; la pressione di uscita all'ottimo trovato
coincide con l'ambiente posto a meno di 916~Pa, dentro una banda
derivata di 7958~Pa; il misuratore di spigolo legge lì
indipendentemente $\approx 0$, quindi la condizione di ottimalità
della teoria di Rao è soddisfatta su un progetto raggiunto dalla
ricerca senza che essa l'abbia mai valutata; e l'obiettivo corrotto
corre fino al bordo dell'intervallo --- rigettato.

\begin{figure}[htbp]
\centering
\includegraphics[width=0.8\textwidth]{fig_jeps.pdf}
\caption{Il paesaggio di ottimizzazione e la fisica dietro al suo
picco, entrambi disegnati dalle 13 valutazioni della marcia
memorizzate dal driver stesso. (a)~Spinta contro forma: rapporto di
espansione $\varepsilon$ [--] sull'asse orizzontale, obiettivo $J$
[MN] su quello verticale. Si osservi quanto è piatta la sommità ---
l'ottimo in forma chiusa (rosso tratteggiato) e il punto di
atterraggio della ricerca alla cieca (blu punteggiato) differiscono
dello 0.09\% in $\varepsilon$ su una curva piatta a parti su
$10^{5}$, quindi la ricerca alla cieca rinuncia a meno di $10^{-5}$
della spinta. (b)~Perché il picco stia lì: pressione di uscita $p_e$
[bar] contro lo stesso $\varepsilon$ [--], in calo attraverso la
linea della pressione ambiente $p_a$. L'attraversamento --- gas che
esce sopra l'ambiente a sinistra (sotto-espanso, allungare), sotto
l'ambiente a destra (sovra-espanso, accorciare) --- cade sullo
stesso $\varepsilon$ del picco in (a). È il lemma dell'area di
uscita, misurato anziché assunto.}
\label{fig:jeps}
\end{figure}
